\documentclass[11pt]{article}
% Uses the canonical latex-docs semantic style in-repository. The local
% lightweight fallback is for independent compilation of this distribution.
\IfFileExists{technical-user-manual.sty}{\usepackage{technical-user-manual}}{\usepackage{../support/aaoa-standalone}}
\usepackage{mathtools}
\setDocTitle{Chapter 8: Strings and Tries}
\setDocSubtitle{Guided reading, derivations and original exercises}
\setDocVersion{1.0}
\setDocStatus{Educational reference / unapproved working material}
\setDocOwner{Algorithm Engineering}
\setDocAuthor{Original companion material}
\setDocSummary{Coverage of primary book Chapter 8 with independently authored examples, proof obligations and lab transfer.}
\providecommand{\Expect}{\mathbb{E}}
\providecommand{\Prob}{\mathbb{P}}
\providecommand{\Var}{\operatorname{Var}}
\providecommand{\bigoh}{\mathcal{O}}
\begin{document}
\makedoctocpage
\section{Reading map and prerequisite check}
\noindent\textbf{Primary source sections:} 8.1 String Searching; 8.2 Bitstrings; 8.3 Regular Expressions; 8.4 Finite-State Automata and KMP; 8.5 Context-Free Grammars; 8.6 Tries; 8.7 Trie Algorithms; 8.8 Combinatorial Properties; 8.9 Larger Alphabets (printed pp. 416--465).\par
\noindent\textbf{Prerequisites:} Completed preceding mathematical-method chapters or equivalent demonstrated mastery.\par
\section{Learning outcomes}
\begin{enumerate}
\item Model search cost in characters, transitions and memory.
\item Demonstrate the role of prefix/suffix reuse in KMP.
\item Compare trie node counts under uniform versus structured key models.
\end{enumerate}
\section{Guided technical analysis}
\subsection{String model and unit}
Specify alphabet $\Sigma$, length $L$, pattern length $m$, distribution of text and pattern, and whether matches overlap. Comparing full strings is different from inspecting individual characters.
\subsection{Automaton perspective}
KMP precomputes the longest proper prefix that is also a suffix for each pattern prefix; on mismatch, it transitions to a feasible state rather than restarting the text scan. This supports linear deterministic search time in text plus pattern length under standard array access assumptions.
\subsection{Trie size}
The number of allocated nodes depends on shared prefixes. For keys of length $L$, an uncompressed character trie has at most $1+mL$ nodes for $m$ inserted keys, but this upper bound can be very loose. A prefix-heavy workload is not interchangeable with uniformly sampled independent strings.
\section{Independent problem set}
\begin{enumerate}
\item Compute failure-function values for a repetitive example pattern and explain each fallback.
\item Construct keysets with many shared prefixes and almost none; compare node allocation and lookup character visits.
\item Explain how a Unicode code-point alphabet differs from bytes and how that impacts the cost unit.
\end{enumerate}
\section{Laboratory transfer}
\textbf{Required laboratory:} LAB-PY04: deterministic trie node counts and prefix-search operations.\par
\noindent\textbf{Minimum recorded evidence:} repository commit or source hash; test inputs/seeds; actual operation counts; predicted values; one explicit model limitation.\par
\section{Mastery gate and self-review}
All lookups are validated; the same corpus and normalization rule are used for every trie variant.\par
\begin{enumerate}\item Can I define the random object/input exactly? \item Is my identity exact, asymptotic, or empirical? \item Which base cases and counterexamples did I check? \item What would falsify my performance hypothesis?\end{enumerate}
\section{Suggested notebook structure}
Use five named sections: model and invariants; derivation; deterministic small instances; randomized replication; discrepancy and revised hypothesis. Keep a proof separate from measurement results.

\end{document}
